\section{Gait Recognition-based Authentication}
\label{sec:new-ideas-gait-recognition}

\paragraph{Chapter \ref{chap:context} Recap} From a security perspective, \ac{gr}-based authentication can contribute to eliminating vulnerabilities found in traditional systems by using inherence factor-based (biometric) credentials tied to user physiology --- similar to existing technologies like Face and Fingerprint recognition (see Section \ref{subsec:biometric-authentication-trends}).
Biometrics remove the need for users to carry physical credentials or remember secret values (e.g., \ac{rfid} Key Cards or \ac{pin} codes in traditional systems).
These approaches reduce the risk of the theft of, or users misplacing or forgetting credentials.

While use of biometric authentication for \ac{pacs} is rising, with implementations of Face, Fingerprint, and Iris recognition being reported \parencite{ifsecglobalStatePhysicalAccess2024}, alternate biometrics like \ac{gr} may help improve the user experience by eliminating pain points experienced with existing solutions.
The most common \ac{pacs} biometrics are inarguably accurate and reliable \parencite{nationalinstituteofstandardsandtechnologyNISTStudyShows2004,universityofbaghdadHighaccuracyModelsIris2024}.
However, as discussed in Chapter \ref{chap:context}, all operate in ways that could be considered obtrusive; users must divert their attention and actively participate (e.g., to scan their fingerprint), or influence accuracy (e.g., face coverings \parencite{cooperEffectsFaceCoverings2022} or changes in environment \parencite{brajeIlluminationEffectsFace1998} affecting Facial Recognition effectiveness).

This project aims to present \ac{cv}-based \ac{gr} as a viable alternative for \ac{pacs} authentication, or authentication in general.
It contrasts existing solutions, by primarily being able to operate passively and identify individuals from a distance, while also not suffering from single points of failure, such as face coverings that impact Face Recognition systems.

Although a promising solution to improving the security and usability of \ac{pacs}, \ac{cv}-based \ac{gr} has its own requirements that may sway effectiveness.
Physically, there must be enough space to enable a camera to capture subject's walking cycle.
It also remains somewhat reliant on environmental conditions, with changes affecting walking style, or camera's ability to capture usable data.
Additionally, as mentioned in Section \ref{subsec:cv-based-gr-techniques}, \ac{gr} requires hardware powerful enough to process gait cycle data in real-time, which is important when trying to identify subjects unobtrusively.

\subsection{Strength, Weakness, Opportunity, and Threat Analysis}

To evaluate this proposal from a strategic perspectice, a SWOT (Strengths, Weaknesses, Opportunities, Threats) analysis has been conducted.
This provides a view of internal and external factors that could influence the success or openness to adoption of \ac{gr}-based authentication.
The proposal is assessed in respect to performance, security, cost, and ethical concerns discussed in Chapter \ref{chap:context}.

\begin{itemize}
    \item \textbf{Strengths}: A \ac{cv}-based \ac{gr} authentication system can operate passively, identifying subjects from a distance without their direct participation.
          It's use of biometric physiological factors (gait), reduces risks relating to lost or stolen credentials.
          A modular client-server architecture supports future scalability and allows an agile development process.

    \item \textbf{Weaknesses}: A \ac{cv}-based \ac{gr} system remains influenced by environmental conditions (e.g., lighting and background contrast) that can affect accuracy, similar to facial recognition.
          Performance of the system may vary on an individual level due to natural gait variability or temporary changes due to injury, or a change in clothing or footwear.
          The real-time demands of a \ac{pacs} may exceed the capabilities of accessible hardware when performing \ac{gr} processing.

    \item \textbf{Opportunities}:
          \ac{gr}-based authentication could be deployed to enable frictionless access control in environments like offices, hospitals, and schools.
          Interest in contactless authentication relates mostly to convenience, security, and hygiene concerns.
          Future developments could combine \ac{gr} with other biometric factors (multi-modal biometrics) for higher accuracy.

    \item \textbf{Threats}:
          Privacy concerns appear high, with primary research (see Section \ref{subsec:perceived-public-opinion-via-survey}) highlighting anxiety about the collection and use of their biometric data.
          The accuracy of the system may prove to be much lower than other biometrics like facial or fingerprint recognition \parencite{nationalinstituteofstandardsandtechnologyNISTEvaluatesFace2021,nationalinstituteofstandardsandtechnologyNISTStudyShows2004}, impacting feasibility.
          Spoofing or adversarial attacks against gait recognition may be easier to achieve compared to other biometrics, which could lead to false acceptances or denials, impacting confidentiality, integrity, or availability.

\end{itemize}
